%% Prime-filter separation and bounded distributive-lattice duality.

\subsection{Prime filters of bounded distributive lattices}
\label{sec:p10-prime-filters}

\begin{hbdefinition}{def:p10-distributive-prime-filter}{Distributive-lattice filters and prime filters}
A bounded distributive lattice \(D\) has finite meet and join
operations, least element \(0\), and greatest element \(1\).
The one-element lattice is permitted.
A filter contains \(1\), is upward closed, and is closed under
binary meets. An ideal contains \(0\), is downward closed, and
is closed under binary joins.
A filter is proper when it excludes \(0\).
A \emph{prime filter} is a proper filter \(P\) such that
\[
 a\vee b\in P\quad\Longrightarrow\quad a\in P\text{ or }b\in P.
\]
A bounded lattice homomorphism preserves \(0,1,\wedge,\vee\).
The principal filter and ideal are
\(\uparrow a=\{d:a\leq d\}\) and
\(\downarrow b=\{d:d\leq b\}\).
The characteristic map of a prime filter is a bounded lattice
homomorphism \(D\to\{0,1\}\), and every such homomorphism arises
this way. This assertion concerns lattice operations.
Preservation of any additional implication operation is a separate
condition.
\index{prime filter!distributive lattice}
\end{hbdefinition}

\begin{hbtheorem}{thm:p10-prime-filter-separation}{proved=annex:p10.priestley.1}{Prime-filter separation}
Assume BPIT. If \(F\) is a filter and \(I\) an ideal of \(D\)
with \(F\cap I=\varnothing\), there is a prime filter
\(P\supseteq F\) disjoint from \(I\).
In particular, if \(a\not\leq b\), a prime filter contains \(a\)
and omits \(b\).
If \(D\) is countable, these conclusions hold in \(\mathsf{ZF}\)
without BPIT.
\end{hbtheorem}

The proof first handles finite distributive lattices, then reduces
the general case to arbitrary-variable propositional compactness.
For the countable case it supplies a separate enumeration construction.
The prime-filter theorem and its Heyting application are treated in
\cite[Prime Filter Theorem and Corollary~2.14, p.~41;
Theorem~2.26, pp.~47--48]{Moraschini2023AlgebraLogicNotes}.

\subsection{Priestley spaces and representation}
\label{sec:p10-priestley-duality}

\begin{hbdefinition}{def:p10-priestley-space}{Priestley space and clopen upsets}
For a partially ordered set \(X\), an \emph{upset} \(U\) satisfies
\(x\in U,\ x\leq y\Rightarrow y\in U\).
A Priestley space is a compact Hausdorff space \(X\) with a partial
order such that
\[
 x\not\leq y\quad\Longrightarrow\quad
 \text{some clopen upset }U\text{ contains }x\text{ and omits }y.
\]
The empty ordered space is allowed.
Write \(\operatorname{ClopUp}(X)\) for its clopen upsets, a bounded
distributive lattice under intersection and union, with bounds
\(\varnothing,X\).
For a bounded distributive lattice \(D\), let
\[
 X_D=\{P:P\text{ is a prime filter of }D\},\qquad
 \widehat a=\{P\in X_D:a\in P\}.
\]
Order \(X_D\) by inclusion and give it the topology generated by
the sets \(\widehat a\) and their complements.
Its basic sets can be written
\(\widehat a\setminus\widehat b\):
finite positive memberships combine by meet and finite negative
memberships by join, with empty meet \(1\) and empty join \(0\).
\index{Priestley space}\index{Priestley duality}\index{clopen upset}
\end{hbdefinition}

\begin{hbtheorem}{thm:p10-priestley-representation}{proved=annex:p10.priestley.2}{Priestley representation}
Assuming BPIT, \(X_D\) is a Priestley space, and
\[
 \eta_D:D\longrightarrow\operatorname{ClopUp}(X_D),
 \qquad \eta_D(a)=\widehat a,
\]
is a bounded lattice isomorphism.
It preserves and reflects the lattice order.
The trivial lattice corresponds to the empty Priestley space.
\end{hbtheorem}

\begin{hbtheorem}{thm:p10-priestley-reconstruction}{proved=annex:p10.priestley.3}{Recovering a Priestley space}
For a Priestley space \(X\), the evaluation map
\[
 \varepsilon_X:X\longrightarrow
      X_{\operatorname{ClopUp}(X)},\qquad
 \varepsilon_X(x)=\{U\in\operatorname{ClopUp}(X):x\in U\},
\]
is an order isomorphism and a homeomorphism.
In particular, every prime filter of its clopen-upset lattice
is the membership filter of exactly one point.
The assertion includes \(X=\varnothing\).
\end{hbtheorem}

\begin{hbtheorem}{thm:p10-priestley-functors}{proved=annex:p10.priestley.4}{Functorial Priestley duality}
Assume BPIT.
Bounded distributive lattices with bounded lattice homomorphisms
are contravariantly equivalent to Priestley spaces with continuous
order-preserving maps.
A homomorphism \(h:D\to E\) gives the map
\[
 X(h):X_E\to X_D,\qquad P\mapsto h^{-1}[P].
\]
A continuous order-preserving \(g:X\to Y\) gives the lattice
homomorphism
\[
 \operatorname{ClopUp}(g):
 \operatorname{ClopUp}(Y)\to\operatorname{ClopUp}(X),
 \qquad U\mapsto g^{-1}[U].
\]
The maps \(\eta_D\) and \(\varepsilon_X\) are natural, and the
two reconstruction composites are identities.
\end{hbtheorem}

These statements give the ordered-space representation of
\textcite{Priestley1970}, with full proofs of the maps and
reconstruction supplied here.
Priestley's later author presentation identifies the hom-functors
and evaluation maps explicitly
\cite[slides~17--20]{Priestley2019Slides}.
The Heyting specialization and its order-topological conventions
are also described in
\cite[\S2.3.3, pp.~33--37]{Bezhanishvili2006Thesis}.

\paragraph{Why the order is part of the representation.}
The three-element chain \(D=\{0<a<1\}\) has exactly two prime
filters, \(P=\{1\}\) and \(Q=\{a,1\}\), with \(P\subsetneq Q\).
The topology is discrete: \(\widehat a=\{Q\}\) and its complement
is \(\{P\}\).
There are exactly three clopen upsets,
\(\varnothing,\{Q\},\{P,Q\}\), reproducing \(D\).
The underlying discrete space has four clopens, including
\(\{P\}\), which is not an upset.
The clopen-upset restriction records the original non-Boolean
lattice structure.
